\section{Discussion}\label{sec:disc}

The engine is far stronger than the humans it met and clearly weaker than Scan.
This section explains how both can be true, why the engine's games against other
programs are almost always drawn, what this means for its rating, what learning
from losses has and has not achieved, and where the results should be read with
care.

\subsection{Why humans lose}\label{sec:disc-humans}

The engine does not beat humans by out-thinking them in the opening. While more
than 30 pieces are on the board the two sides play about equally well
(\cref{fig:errors}). The difference is consistency. In the middlegame and
endgame the humans gave away three to five times as much per move as the engine,
and every so often lost a man or more in a single move, which the engine did not
do once in more than three thousand moves. A human playing blitz has a few
seconds per move and will sooner or later overlook a capture sequence; the
engine checks every capture sequence on every move. Once a game was decided, the
engine converted it: no human game it was winning by a man or more ended in a
draw. Its results against humans correspond to a performance rating of about
2640 (a 96.7\% score against opponents averaging 2049, by \cref{eq:elo}), almost
200 points above its actual rating.

\subsection{Why engines draw}\label{sec:disc-draws}

Against bots the picture reverses. Neither side made errors large enough to
lose, and 44 of 51 games were drawn. This is a property of the game as much as of
the players. The $8\times8$ form of draughts has been proven a draw with perfect
play \parencite{schaeffer2007}. International draughts has not been solved, but
it shows the same tendency: even in our benchmark against Scan, an opponent about
93 Elo stronger, 47\% of games were drawn, and when Scan and valhalla played out
12 positions in which valhalla had judged itself a man ahead, each taking the
better side in turn, not one of the 48 games was won. A small advantage in draughts is usually not
enough. With fewer pieces the defending side can often reach one of the endgames
in which the draw rules of \cref{sec:rules} end the game after 5 or 16 moves, and
most of the drawn games ended exactly that way (\cref{tab:endings}). Between
strong programs, a stronger engine therefore turns into a higher score only
slowly: it wins slightly more often, but it mostly draws.

\subsection{What limits the rating}\label{sec:disc-rating}

A Glicko-2 rating settles where the score a player is expected to make matches the
score it actually makes. Against an opponent rated 400 points lower the expected
score is about 91\%, and against one rated 600 points lower about 97\%. In
practice this made draws expensive and wins cheap. Once its rating was
established, a win against a human rated around 2050 earned the engine between 1
and 8 points, and by the end only 1 or 2. Each of the three draws against humans
rated about 300 to 460 points lower cost it between 8 and 28 points. A
single such draw undid several wins. Draws against bots rated
around 2400 held the rating near that level. The rating is therefore limited less
by how the engine plays than by who it plays: to climb further it must score
nearly 100\% against humans, draw almost never, and win some of its games against
the other bots, which, as \cref{sec:disc-draws} explains, is difficult in a game
this drawish. A rating far above the best human players would require opponents
who are rated far above them as well, and the site does not have them.

\subsection{Learning from losses}\label{sec:disc-learning}

Because the engine never lost to a human, the lesson procedure of
\cref{sec:learning} ran only for draws against much weaker humans (an extension
we added during the experiment) and for losses in the sparring rounds against
Scan. Its effect on strength is small and hard to measure. A 200-game match can
only detect differences of about 30 Elo or more, so the acceptance rule of 49\%
lets through networks that are no better than the one they replace: champ\_005
scored within 3 Elo of champ\_004 against Scan (\cref{tab:scan}). The sparring
rounds against Scan gave three accepted networks, each scoring between 51.5\%
and 53\% against its predecessor, and then six rounds in a row that were
rejected. The largest gain came instead from training a new network from scratch
on more varied data (v2, 20 Elo against Scan). More data of the same kind did
not help: v3, trained like v2 but with the deeper Scan self-play set grown from
7.3 to 11.1 million positions, was weaker. This suggests that what the network lacks is not more examples, but
examples from the positions it actually reaches in its own games, which Scan's
self-play does not cover. The two losses on the board point the same way: the
engine was not caught by a single oversight but slowly outplayed in the
middlegame, which suggests a weakness in its judgement of quiet positions rather
than in its calculation.

\subsection{Limitations}\label{sec:disc-limits}

The results cover three days and 127 games. The 76 games against humans were
played against 16 players, two of whom account for 46 games, and the strongest
of them was rated 2293. The highest-rated humans we challenged did not accept, so the claim that no human has beaten the engine holds for players
rated up to about 2300. All games were blitz, where humans make more mistakes
than with more time to think; the gap between engine and humans is likely smaller
at slower time controls. The opponents chose to play a bot, which may not make
them representative of the site's players. The analysis of how games were
decided relies on Scan, which is also the engine's teacher, so the two may share
blind spots; its scores carry noise from a finite search, and the thresholds we
used (one man for a decided game, 30 and 15 pieces between opening, middlegame
and endgame) are our own choices. Finally, three of the five losses were caused
by a fault in the engine's time management rather than by its play; they say
little about its strength but show that a deployed engine must be tested against
the conditions of the server, not only against itself.
